\honest{Changing Your Metaethics}{Eliezer Yudkowsky}{2008}

If you say ``Killing people is wrong,'' that is morality. If you say you should not kill because God forbade it, that is metaethics. People agree much more about the first than about the second.

Theists, Objectivists and Platonists all agree that killing is wrong, disagree about what ``wrong'' means, and insist that without their own metaethic morality falls apart. Most of them must be wrong, so many people will have to change their metaethics. I have learned that this is ``nigh-impossible in the presence of an unanswered attachment,'' so I offer lines of retreat, collected from my earlier posts. \nb{Seven of the posts summarized or linked here are not in this book: ``The Moral Void,'' ``Is Morality Given?'', ``Rebelling Within Nature,'' ``Does Your Morality Care What You Think?'', ``Math is Subjunctively Objective,'' ``Probability is Subjectively Objective'' and ``Existential Angst Factory.''}

To someone who thinks morality comes from God: ``Take personal responsibility!'' Even if God gave you orders, it would be your own decision to obey them. The same goes for any source of moral authority. \nb{Kant made a closely related point: even the Holy One of the Gospels must first be compared with our ideal of moral perfection. In the dialogue the post links to, ``Is Morality Given?'', Obert grants the point and adds that it does not follow that morality consists merely of one's preferences. The post uses the argument only to reassure, which stays within that limit.}

The most important line of retreat: if your metaethic stops telling you to save lives, you can drag the kid off the train tracks anyway. ``Maybe that which you would do even if there were no morality, is your morality.'' My point is not ``that no morality exists.'' \nb{Even error theorists, who hold that no moral claim is true, can keep moral language in a new sense or replace it with another way of speaking.}

Then I summarize other posts. You cannot explain moral arguments to a rock. My two posts on reflection ``explain the difference'' between a reflective loop and circular logic. Evolution's causal role in morality is not a justifying role. And a quantity computed inside your head need not be about your thoughts: a calculator can compute ``What is 2 + 3?'' rather than its own output. \nb{The worry this answers, that whatever you think is right must then be right, is a standard objection to simple subjectivism.}

I promise to present my own metaethic in the next posts. \nb{Those posts, ``Setting Up Metaethics'' and ``The Meaning of Right,'' are not in this book. The book places this post before ``Could Anything Be Right?'', which was written nine days earlier.}

Metaethics matters, although it adds up to moral normality, because many people are harmed by misunderstanding where their morality comes from. ``Poor metaethics forms part of the teachings of many a cult, including the big ones.'' I name no such teaching. Then I start to give ``the real reason'' and stop.
